\chapter{Performance}

Two expectations in the original design were contradicted by measurement, and the first
explanation of one of them was then contradicted as well. They are written down here
because an expectation that was wrong and got corrected is more useful than one that was
never recorded.

\entry{Kernel fusion, expected ten to twenty percent, measured about one}

\field{The expectation.} Fusing the force evaluation with the Verlet update takes one launch
instead of three and stops reading the acceleration array. Ten to twenty percent seemed
about right for a memory bound kernel.

\field{The measurement.} Round by round, at the default tile width, fused over tiled is
\FusedOverTiledRatioAtTenThousand{} at \LadderMinN{} particles and
\FusedOverTiledRatioAtHundredFiftyThousand{} at one hundred and fifty thousand, both beyond
their interquartile ranges and the session drift of \RepeatabilitySpread{} percent. At
\LadderMaxN{} particles it is \FusedOverTiledRatioAtMaxN{} with a range of
\FusedOverTiledRatioIqrAtMaxN{}, which does not resolve anything: the tiled trials are
bimodal, near \TiledLowModeMsAtMaxN{} and \TiledHighModeMsAtMaxN{} milliseconds. I do not
yet know why.

\field{Root cause of the wrong expectation.} Arithmetic intensity. The all pairs sum performs
$O(N^2)$ arithmetic on $O(N)$ data, so the ratio of operations to bytes grows without limit
with the problem size, and device memory is nearly idle: Nsight Compute puts it at
\NaiveDramPct{} percent of peak for the naive kernel and \TiledDramPct{} percent for the
tiled one.

\field{What I got wrong next.} \texttt{v1.0.0} turned this into a ceiling. It called the saving
a round trip of the acceleration array, $24N$ bytes written and read, priced that at the
measured bandwidth, and said no implementation could beat the result. Three things were wrong
with that. The fused kernel still wrote the array. The split path reads the array twice,
once in each integrator launch, so the description of the traffic was wrong too; counted from
the code, fusion removes \ModelBytesPerParticleSaved{}$N$ bytes, which is
\ModelSavingMsAtMaxN{} milliseconds or \ModelSavingPercentAtMaxN{} percent of a step at
\LadderMaxN{} particles. And the measured gains are several times larger than any such
figure, so it was never a ceiling. It prices one resource and says nothing about the rest of
the step. The gain itself, 0.1 percent in \texttt{v1.0.0}, was a difference of 0.12 percent
between separate blocks against a session spread of about 3 percent, so it could not be told
from noise.

\field{What I changed.} The protocol, so a difference of a percent can be read. The source of
the gain is still unmeasured. The fused kernel stays because it is faster where the difference
is resolved; \texttt{v1.0.0} also said it halves the memory footprint of a step, which is false,
because it adds a second position buffer.

\entry{Shared memory tiling is not a bandwidth optimisation on this problem}

\field{The expectation.} Tiling is conventionally introduced as a way of
reducing device memory traffic, and I introduced it that way. The prediction
that follows is that the naive kernel is limited by memory bandwidth and the
tiled kernel relieves it.

\field{The measurement.} Tiling delivers a factor of \TilingSpeedupAtMaxN{},
which is real and large. But device memory throughput is
\NaiveDramPct{} percent of peak before the change and \TiledDramPct{} percent
after. Device memory was never the constraint, so the conventional explanation
of why tiling helps is wrong here even though the tiling itself is right.

\field{What is actually happening.} The counters point at the first level cache
and the load store pipeline:

\begin{center}
\begin{tabular}{l rrr rr}
\toprule
& \multicolumn{3}{c}{throughput, percent of peak} & \multicolumn{2}{c}{per launch} \\
\cmidrule(lr){2-4}\cmidrule(lr){5-6}
variant & L1 and load store & L2 & device memory & global load sectors & warp issue \\
\midrule
naive & 87.98 & 1.62 & 0.34 & 5,000.1M & 54.99 \\
tiled & 45.41 & 1.61 & 0.35 & 78.3M & 82.68 \\
fused & 45.90 & 1.64 & 0.31 & 78.3M & 83.59 \\
\bottomrule
\end{tabular}

\end{center}

Global load sectors fall by a factor of \SectorReduction{} and load store
instructions by a factor of \LsuReduction{}. First level cache throughput falls
from \NaiveLoneTexPct{} percent of peak, close to saturation, to
\TiledLoneTexPct{} percent. The decisive number is warp issue efficiency, which
rises from \NaiveIssuePct{} percent to \TiledIssuePct{} percent. Tiling moves the reads into
shared memory and the load store pipe stops being the limit.

\field{Why this mattered practically.} Attacking the naive kernel's device memory traffic would
not have addressed the unit the counters name as limiting it. The profile pointed at the right
resource; reasoning from the familiar description would not have.

\field{Control.} The arithmetic instruction counts are equal across all three
variants to within a fraction of a percent, which is what licenses attributing
the difference to data movement rather than to the variants computing different
things. Those same counters also let the conventional
\ConventionalFlopsPerInteraction{} floating point operations per interaction be checked
against the \MeasuredFlopsPerInteraction{} actually retired, counting a fused multiply add as
two. Interactions per second is the measurement, and TFLOP/s depends on which count is used.

\entry{The first default tile width is not the fastest}

\field{Symptom.} A sweep of tile widths at \TileSweepParticles{} particles shows a width of
thirty two, the \texttt{v1.0.0} default, is well behind the widest of the four measured.

\begin{center}
\begin{tabular}{r rr rr}
\toprule
& \multicolumn{2}{c}{tiled} & \multicolumn{2}{c}{fused} \\
\cmidrule(lr){2-3}\cmidrule(lr){4-5}
tile width & ms per step & TFLOP/s & ms per step & TFLOP/s \\
\midrule
32 & 12.116 & 16.51 & 12.110 & 16.52 \\
64 & 11.777 & 16.98 & 11.769 & 16.99 \\
128 & 11.585 & 17.26 & 11.490 & 17.41 \\
256 & 11.487 & 17.41 & 11.373 & 17.59 \\
\bottomrule
\end{tabular}

\end{center}

\field{Root cause.} Two barriers are paid per tile, so a narrower tile pays them more often per
unit of arithmetic. That fits the direction of the result, though the barrier cost was not
measured on its own. A width of \BestTileWidth{} runs at \BestTileMs{} milliseconds against
\TileThirtyTwoMs{}, a difference of \TileWidthGainPercent{} percent, and a paired ratio of
\BestTileOverTileThirtyTwoRatio{} with an interquartile range of
\BestTileOverTileThirtyTwoRatioIqr{}.

\field{Fix and why.} \texttt{v1.0.0} kept thirty two and reported the sweep. The default is now
\BestTileWidth{}, and the ladder was measured there, because the difference is far outside the
session drift of \RepeatabilitySpread{} percent.

\field{A defect in the sweep itself.} The \texttt{v1.0.0} sweep timed the force kernel alone for
tiled, with particles that never moved, against whole steps for fused, so any comparison across
variants in its table was void. It also carried a claim here that thirty two was the slowest of
the four, when its own table had sixty four slower. Both variants now time a whole step.

\entry{Absolute timings depend on the thermal state of the card}

\field{Symptom.} In \texttt{v1.0.0}, a tile width sweep run immediately after the full benchmark
sweep produced timings about fifteen percent higher than the same sweep run from cold. Nothing
in the code had changed.

\field{Root cause.} Sustained clocks falling as the card heats is the likely reason, though I did
not log clocks. Either way it is large enough to swamp several of the effects being measured
here, including the entire fusion result.

\field{First fix, and why it was not enough.} \texttt{v1.0.0} compared within one session, kept the
fastest of three trials, and claimed relative differences were stable to a tenth of a percentage
point. They were not. Later measurement of the same binaries found them 4.7 percent apart in
separate blocks and 1.5 percent apart the other way when interleaved, because drift is confounded
with which variant ran when, and more trials do not help.

\field{Fix and why.} Measure in rounds: every configuration's first trial, then every
configuration's second, with the order rotating. Report the median of ten trials with its
dispersion. Take comparisons from the ratio inside each round, because two configurations in one
round saw the same machine. Measure one reference configuration at the start, middle and end of
the sweep, and treat the spread of those, \RepeatabilitySpread{} percent here, as the smallest
difference worth quoting.

\field{Verification.} Rerunning the sweep on a later commit reproduced the tile ratio to within a
tenth of a percent, and \texttt{scripts/check\_provenance.py} holds every recorded row to a commit
that resolves and a clean tree.
